% Reciprocal proofs for the Mac Lane--Moerdijk source-specific semantics.
\section{Mac Lane--Moerdijk internal semantics}

\begin{annexitem}{annex:p10.mlm.1}{thm:p10-mlm-presheaf-predicates}{Presheaf residuals, quantifiers and base change}
\begin{proof}
Arbitrary componentwise unions and intersections of subpresheaves are closed
under restriction; they have their asserted upper- and lower-bound properties
by component containment. Finite intersection distributes over unions because
the corresponding set operations do. The empty union is the empty presheaf,
and the empty intersection inside \(X\) is \(X\).

The displayed implication is restriction-stable: if \(x\) satisfies its test
at \(c\) and \(u:d\to c\), then every further \(v:e\to d\) tests \(x\cdot(uv)\),
so \(x\cdot u\) satisfies the same condition at \(d\).
If \(R\subseteq A\Rightarrow B\), an \(x\in R(c)\cap A(c)\) satisfies the
implication test with \(1_c\), hence belongs to \(B(c)\).
Conversely, if \(R\cap A\subseteq B\), \(x\in R(c)\) and \(x\cdot u\in A(d)\),
restriction stability gives \(x\cdot u\in R(d)\cap A(d)\subseteq B(d)\).
This proves both implications of residuation. Negation is obtained by setting
\(B=0\). In particular the displayed operations form a Heyting algebra,
not merely a distributive lattice with a named binary operation.

For \(f:X\to Y\), inverse image is \((f^*B)(c)=f_c^{-1}[B(c)]\).
It is restriction-stable by naturality. Ordinary inverse images preserve
unions, intersections, empty and full subsets. For implication, membership
of \(x\in X(c)\) in \(f^*(B\Rightarrow D)\) requires, for every \(u:d\to c\),
\(f_c(x)\cdot u\in B(d)\Rightarrow f_c(x)\cdot u\in D(d)\).
Naturality replaces \(f_c(x)\cdot u\) with \(f_d(x\cdot u)\). This is exactly
the test for \(x\in(f^*B\Rightarrow f^*D)(c)\). Negation follows too.

The existential image is a subpresheaf: a witness \(x\in A(c)\) restricts
to \(x\cdot u\in A(d)\) and \(f_d(x\cdot u)=f_c(x)\cdot u\).
Its containment in \(B\) is equivalent to every \(f_c(x)\), \(x\in A(c)\),
belonging to \(B(c)\), which is exactly \(A\subseteq f^*B\).
For the universal image, restriction stability follows by replacing each
\(v:e\to d\) after restricting along \(u:d\to c\) by \(uv:e\to c\).
If \(B\subseteq\forall_f A\), any \(x\in(f^*B)(c)\) is tested by \(1_c\)
and lies in \(A(c)\). Conversely if \(f^*B\subseteq A\), \(y\in B(c)\) and
\(f_d(x)=y\cdot u\), then \(y\cdot u\in B(d)\), so \(x\in A(d)\).
Thus \(f^*B\subseteq A\) iff \(B\subseteq\forall_f A\).
The restrictions on the target value in this proof are essential for typing.

For Frobenius, \(y\in\exists_f(A\cap f^*B)(c)\) says there is an \(x\in A(c)\)
with \(f_c(x)=y\) and \(f_c(x)\in B(c)\). This is exactly
\(y\in(\exists_f A)(c)\cap B(c)\).
In the pullback, use the canonical component
\(P(c)=\{(x,z):f_c(x)=g_c(z)\}\), with coordinate projections and restrictions.
For \(z\in Z(c)\), membership in \(g^*\exists_f A\) is the existence of
\(x\in A(c)\) with \(f_c(x)=g_c(z)\), exactly the image condition for
\(\exists_qp^*A\). Membership in \(g^*\forall_f A\) is
\[
 \forall u:d\to c\ \forall x\in X(d),\quad
 f_d(x)=g_d(z\cdot u)\ \Longrightarrow\ x\in A(d).
\]
Membership in \(\forall_qp^*A\) tests all \(u:d\to c\) and all pairs
\((x,z\cdot u)\in P(d)\), asking that their first coordinate lie in \(A(d)\).
It is precisely the same condition. This proves both base-change identities
without asserting them for a merely commuting nonpullback square.
All arguments quantify only over existing elements; none selects a family of
witnesses. Empty component sets and the empty base therefore satisfy the
same proofs.
\end{proof}
\end{annexitem}

\begin{annexitem}{annex:p10.mlm.2}{thm:p10-mlm-internal-representation}{Internal representation without a global-element shortcut}
\begin{proof}
The power-object correspondence is the actual inverse bijection of
Theorem~\ref{thm:p10-presheaf-power-objects}, natural in its parameter.
Each component of \(P(X)\) is \(\operatorname{Sub}(\mathrm yc\times X)\).
Restriction along \(u:d\to c\) is inverse image along \(\mathrm yu\times1_X\),
which preserves Heyting operations by the preceding proof. Hence these
component operations are natural maps. Every equational Heyting identity
holds at every component, so is an identity of natural maps. The order
equalizer gives ordinary subobject containment at every component.

For quantifier naturality, the square formed by \(1\times f\) horizontally
and \(\mathrm yu\times1\) vertically is a pullback: it is so at every set
component by its product-coordinate description. Both Beck--Chevalley
identities prove that \(\exists_{1\times f}\) and \(\forall_{1\times f}\) commute
with restriction. The component adjunctions therefore give the two internal
order adjunctions. Explicitly the inequalities that are their unit and
counit hold at each component and thus factor through the internal order
equalizers. Those unit--counit inequalities imply the order adjunction on
every hom-poset by composing with a parameter arrow, so this is internal
adjunction, not only an adjunction on global elements.

Let \(k:Z\to P(X)\) classify \(S\subseteq Z\times X\). For \(z\in Z(c)\),
write \(\bar z:\mathrm yc\to Z\), \(\bar z_d(u)=z\cdot u\).
Then \(k_c(z)\) classifies \((\bar z\times1)^*S\).
For a Heyting operation, the already proved preservation under this inverse
image identifies the operation on \(k_c(z)\) with the classifying element of
the external operation on \(S\). For either quantifier, apply Beck--Chevalley
to the pullback with vertical arrows \(\bar z\times1\). It identifies the
quantifier of \(k_c(z)\) with the classifying element of
\(\exists_{1_Z\times f}S\) or \(\forall_{1_Z\times f}S\).
This proves the claimed action for every \(Z\), element and component.
It also identifies \(P(f)\) with external inverse image.

Apply external Frobenius or either base-change identity after multiplying
the relevant square by an arbitrary parameter \(Z\). Products preserve its
pullback property componentwise. Transport the resulting identity through
the classification bijection. It is an identity after precomposition by
every parameter map. Taking the parameter to be the domain of the two maps
and its map to be identity proves equality of the maps themselves.
This is the Yoneda/representability proof route: identities are detected by
all generalized elements, not by global elements alone. It also fills the
internal Frobenius and universal Beck--Chevalley checks assigned to Exercise~9
in the source's IV\S9.
\end{proof}
\end{annexitem}

\begin{annexitem}{annex:p10.mlm.3}{thm:p10-mlm-contextual-semantics}{Context, substitution, extensionality and comprehension}
\begin{proof}
The five foundation constructions provide products, exponentials, diagonals,
classifier and power objects; the preceding two theorems provide the internal
Heyting operations and both quantifiers. Thus each inductive language clause
has its stated arrow type. Variable projections, composition and pairing
are arrows, and their equations are their actual universal properties.

For substitution, use induction on terms and formulas. A variable projection
composed with the tuple \(a:H\to G\) is exactly the substituted term.
Unit, arrow application, projection and pairing commute with composition.
Application does too because
\(\langle r,s\rangle a=\langle ra,sa\rangle\).
For lambda abstraction, after renaming the bound variable to avoid capture,
the substituted body is \(t\circ(a\times1_X)\). Naturality of transposition
gives \(\Lambda(t\circ(a\times1_X))=(\Lambda t)\circ a\).
Equality and membership commute with composition of their characteristic
maps, so their extensions pull back. The Heyting connectives commute with
inverse image by Theorem~\ref{thm:p10-mlm-presheaf-predicates}.
For quantification, the context square
\[
\begin{array}{ccc}
 H\times X&\xrightarrow{a\times1_X}&G\times X\\
 \pi_H\downarrow&&\downarrow\pi_G\\
 H&\xrightarrow{a}&G
\end{array}
\]
is a pullback. Both Beck--Chevalley identities give
\(a^*\exists_{\pi_G}[\phi]=\exists_{\pi_H}(a\times1_X)^*[\phi]\)
and the analogous universal equation. This completes the induction,
including formulas nested under lambda or quantifiers. It proves the
formula-arrow identity as well, since the classifier of an inverse image is
the original characteristic map composed with the substitution arrow.

For comprehension, let \(k=\Lambda\phi:G\to P(X)\). The exponential inverse
identity says \(\operatorname{ev}(k\times1_X)=\phi\).
Pulling truth back along this equality gives exactly \([\phi]\), hence
membership in the predicate-valued comprehension has the original extension.
This is the full power-object construction, not a selection of unrelated
component subsets. The same exponential inverse identities give beta
conversion \((\lambda x.t)(s)=t[s/x]\) and eta conversion
\(\lambda x.r(x)=r\) when \(x\) is not free in \(r\).

Equality of two terms is the pullback of the diagonal along their tuple,
so it is their equalizer. On this equalizer the two arrows agree; composing
any arrow with them or applying the substitution identity makes the
resulting expression interpretations agree there. This proves semantic
replacement of equals, including formulas and bound-variable expressions.

For predicate extensionality at \(\gamma\in G(c)\), let the interpreted
predicate elements \(r_c(\gamma),s_c(\gamma)\) be subpresheaves
\(R,S\subseteq\mathrm yc\times X\). Their membership after restriction
\(u:d\to c\), at \(x\in X(d)\), is respectively \((u,x)\in R(d)\)
and \((u,x)\in S(d)\), by power-object classification.
Universal quantification of the two implications forces containment both
ways for every such pair: to see one direction, test implication with
identity after each \(u\); for the converse, equality of \(R,S\) makes each
implication true after every further restriction. Thus the universal
equivalence holds exactly when \(R=S\), precisely component equality of
the two predicate-valued terms. The same argument for function elements
uses \(r_c(\gamma),s_c(\gamma):\mathrm yc\times X\to Y\): quantified
pointwise equality equates their value at every \((u,x)\), and hence all
components of these natural maps. Conversely equal maps have equal
evaluations. This proves function extensionality.

The empty hypothesis intersection is the top subobject, so the consequence
and validity definitions agree when there are no hypotheses.
By classification, \([\phi]=G\) iff its characteristic arrow is truth
composed with \(!_G\). For a projection \(\pi:H\times X\to H\), the universal
adjunction gives \(H\subseteq\forall_\pi A\) iff \(H\times X\subseteq A\).
Therefore \(\forall_\pi A=H\) iff \(A=H\times X\). Repeating this for
the finite declared context proves the universal-closure assertion, even
when a type is empty. This does not assert that dropping an unused variable
preserves validity. No step establishes soundness of an unenumerated
formal calculus.
\end{proof}
\end{annexitem}

\begin{annexitem}{annex:p10.mlm.4}{thm:p10-mlm-generalized-clauses}{Generalized-element clauses with their exact dependencies}
\begin{proof}
Use the earlier elementary-topos results explicitly stated in
Theorem~\ref{thm:p10-mlm-topos-structural-package}.
This is a derivation of the clauses from those structural results, not a
new proof of all the source's IV\S6--9. In presheaves those ingredients have
the concrete proofs linked in that statement.
Write \(A=[\phi]\), \(B=[\psi]\), and \(M=\operatorname{Im}(\alpha)\subseteq G\).

Persistence is composition of a factorization. For descent, let \(e:V\to U\)
be epic and suppose \(\alpha e\) factors through \(A\). Its characteristic
equation is \(\chi_A\alpha e=\mathsf{true}\,!_Ue\). Cancel \(e\) to obtain
\(\chi_A\alpha=\mathsf{true}\,!_U\); the characteristic pullback supplies
the factorization. Top is immediate. Bottom forces \(U\to0\), which by
strictness of \(0\) implies \(U\cong0\); conversely an initial domain
factors through every subobject. Conjunction follows from the intersection
pullback universal property.

For disjunction, \(A\vee B\) is the image of \(A+B\to G\).
If \(\alpha\) factors through it, pull this image's epic map back along that
factorization. Pullback stability of coproducts identifies the resulting
domain with \(V+W\), where \(V=U\times_G A\), \(W=U\times_G B\).
The copairing to \(U\) is epic, and its two composites factor through \(A\),
\(B\). Conversely, if such a copairing is epic, the composite from \(V+W\)
factors through \(A\vee B\). Epimorphic descent proves that \(\alpha\) does.
This is the source's image/coproduct route.

For implication, if \(M\subseteq A\Rightarrow B\), every \(\alpha p\) factoring
through \(A\) has image in \(M\cap A\subseteq B\), so forces the consequent.
Conversely pull \(A\cap M\hookrightarrow M\) back along the epic map
\(U\to M\) to obtain \(p:V\to U\). The induced map \(V\to A\cap M\) is
epic, so \(\operatorname{Im}(\alpha p)=A\cap M\). The antecedent holds
on \(V\); the testing condition therefore gives \(A\cap M\subseteq B\).
Residuation implies \(M\subseteq A\Rightarrow B\).
For negation set \(B=0\) in this proof. A refining domain forces bottom
exactly when it is initial, which proves the omitted negation verification.

Let \(R=[\phi]\hookrightarrow G\times X\), with image \(D=\exists_\pi R\)
under \(\pi:G\times X\to G\). If \(\alpha\) factors through \(D\), pull the
epi \(R\to D\) back along its factorization. This gives an epic \(p:V\to U\)
and a map \(V\to R\); its \(X\)-coordinate is \(b\), and its
\(G\)-coordinate is \(\alpha p\). Thus the required witness exists on \(V\).
Conversely the given witness map factors through \(R\), hence its
\(G\)-coordinate \(\alpha p\) factors through \(D\). Descent proves forcing
on \(U\). No section of \(p\) has been chosen.

For universality, the universal adjunction gives
\(M\subseteq\forall_\pi R\) iff \(M\times X\subseteq R\).
Pullback-stable images give
\(\operatorname{Im}(\alpha\times1_X)=M\times X\).
Thus this is precisely forcing at the generic tuple
\((\alpha\pi_1,\pi_2):U\times X\to G\times X\).
If that tuple factors through \(R\), every \((\alpha p,b)\) factors through
it by \(\langle p,b\rangle:V\to U\times X\). Conversely the all-arrow test
includes \(V=U\times X\), \(p=\pi_1\), \(b=\pi_2\), proving equivalence.

For equality the extension is a pullback of the diagonal, so its
factorization condition is equality of the two interpreted arrows.
Membership is the pullback of truth along evaluation. Pulling this square
back along the element--predicate tuple shows that its forcing condition is
exactly factorization through membership. This expands the source's
classifier/evaluation argument without a global-element assumption.
The same reasoning applies to any finite tuple of variables; all products
are in the topos, not in the base category of a site.
\end{proof}
\end{annexitem}

\begin{annexitem}{annex:p10.mlm.5}{thm:p10-mlm-site-clauses}{The explicit subsheaf route and trivial-topology specialization}
\begin{proof}
We work with the stated sheaf category and topology, without constructing
associated sheaves. For a restriction-stable subpresheaf \(T\) of a sheaf
\(F\), put
\[
 L(T)(c)=\{z\in F(c):\{u:d\to c:z\cdot u\in T(d)\}\in J(c)\}.
\]
The braces form a sieve. Stability of covering sieves makes \(L(T)\)
restriction-stable. It contains \(T\), because a section of \(T\) has the
maximal membership sieve. If \(z\)'s restrictions on a cover belong to
\(L(T)\), then the membership sieve of \(z\) pulls back to a cover along
each arrow of that cover. Transitivity makes it cover, so \(z\in L(T)\).
This proves local membership. A subpresheaf of a sheaf with local membership
is itself a sheaf: amalgamate its matching family in \(F\); the amalgamation
locally belongs and hence belongs to the subpresheaf. Uniqueness holds in
\(F\). Conversely every subsheaf has local membership by its amalgamation
axiom and the uniqueness of amalgamation in \(F\).
Consequently \(L(T)\) is the least subsheaf containing \(T\): every such
subsheaf contains all the locally \(T\)-valued sections by local membership.

Meets are intersections; joins are \(L(A\cup B)\), which gives the disjunction
clause. The least subsheaf is \(L(0)\). Its membership sieve is empty, so
it has a section at \(c\) exactly when \(\varnothing\in J(c)\).
If the empty sieve covers, the sheaf axiom for its unique empty matching
family gives exactly one section of every sheaf at that object, matching
the source's III\S8 (18).

Define \(D(c)\) by the all-arrow implication condition
\(z\cdot u\in A(d)\Rightarrow z\cdot u\in B(d)\).
Composition proves restriction stability. If restrictions of \(z\) on a
cover satisfy it, test \(u:d\to c\) with \(z\cdot u\in A(d)\). Pull that
cover back to \(d\). On every arrow \(v:e\to d\) of its pullback, the
implication gives \(z\cdot uv\in B(e)\). Local membership of \(B\) gives
\(z\cdot u\in B(d)\). Thus \(D\) is a subsheaf.
For any subsheaf \(R\), \(R\subseteq D\) implies \(R\cap A\subseteq B\)
by testing identity. Conversely if \(R\cap A\subseteq B\), restriction
stability of \(R\) gives the all-arrow condition for every element of \(R\).
This proves residuation, hence identifies \(D\) as Heyting implication.
Taking \(B=L(0)\) gives the negation clause.

For a sheaf map \(f:X\to Y\) and subsheaf \(A\subseteq X\), its existential
image is \(L(I)\), where \(I(c)=f_c[A(c)]\). Naturality makes \(I\) a
subpresheaf. For a subsheaf \(B\subseteq Y\),
\(L(I)\subseteq B\) iff \(I\subseteq B\) iff \(A\subseteq f^*B\).
This proves the left adjunction and its covering-witness formula.
There is no family of witnesses chosen across all stages: the sieve asks
for existence independently at each arrow. If a family formulation with
selected witnesses is wanted, index the cover by all pairs of an arrow and
one of its witnesses. Its generated sieve is the same covering sieve and
the witnesses are then tautologically available, without using Choice.

Define the universal image using the same all-arrow formula as for
presheaves. It is restriction-stable by composition. To prove local
membership, suppose \(y\)'s restrictions on a cover \(S\) satisfy the formula.
Take \(u:d\to c\), \(x\in X(d)\) with \(f_d(x)=y\cdot u\).
For every \(v:e\to d\) in \(u^*S\), naturality gives
\(f_e(x\cdot v)=y\cdot uv\); the condition on the corresponding covered
restriction of \(y\) implies \(x\cdot v\in A(e)\).
Local membership of \(A\) gives \(x\in A(d)\). Thus the universal image is
a subsheaf. The two directions of its order adjunction are as in
Annex~\ref{annex:p10.mlm.1}, using restriction stability and identity.
Specialize these images to the projection of the quantified variable to
obtain existence and universality. Equality and membership are pullback
predicates, so their component tests are unchanged. All these operations
have been characterized by their universal properties, which identifies
them with the internal-language extensions, not just candidate subsets.

Restriction persistence follows from subfunctor closure and cover locality
from subsheaf local membership. For the trivial topology, a covering sieve
contains identity. Thus covering-witness existence entails a witness at
identity; a covering-disjunct clause entails one disjunct there. Conversely
a current-stage witness or disjunct restricts to all arrows and supplies
the maximal cover. The empty sieve is not maximal at any object, so bottom
never holds, and negation excludes every positive restriction.
No base-category products are used. If the base is empty there are no
stages to test, so validity is vacuous in its degenerate presheaf topos.

Finally, the site/generalized-element comparison uses precisely the stated
\(a\dashv i\) and Yoneda bijections: a section's membership in a subsheaf
is equivalent to factorization of its map from \(a(\mathrm yc)\).
This is not a proof of the associated-sheaf construction. The alternative
source proof by a generating family of sheafified representables has the
separate prior dependencies identified in the body. For the dense topology
on a poset, substitute its covering criterion into the disjunction and
existence clauses. Covers cannot be empty, since \(q=p\) must have a further
\(r\leq p\); the negation clause consequently becomes absence of every
positive refinement. This gives the final dense-topology illustration.
\end{proof}
\end{annexitem}

\begin{annexitem}{annex:p10.mlm.6}{thm:p10-mlm-internal-epi-boolean-choice}{Epis, Booleanness and internal choice}
\begin{proof}
We use exactly the generalized clauses and the structural package already
identified, including pullback-stability of epimorphisms.
Let \(h=\langle\pi_1,\widetilde r\rangle:U\times X\to U\times Y\).
The generic-element universal clause changes internal surjectivity into
\(U\times Y\forces\exists x:X\,r(\pi_1)(x)=\pi_2\).
The existential clause gives an epi \(e:V\to U\times Y\) and \(b:V\to X\)
such that \(\widetilde r\langle\pi_1e,b\rangle=\pi_2e\).
Consequently \(h\langle\pi_1e,b\rangle=e\); since a map through which an
epi factors is epic, \(h\) is epic. Conversely if \(h\) is epic, choose
\(V=U\times X\), \(e=h\), \(b=\pi_2\), which satisfies the required
equation. This proves the first equivalence without constructing a choice
of inverses. It is also the source's representation of its epi subobject.

Every subobject \(A\hookrightarrow U\) has a characteristic
\(a:U\to\Omega\). The formula consisting of the truth-value variable
itself, after substitution of \(a\), has extension \(A\).
If universal excluded middle holds, the universal clause makes
\(U\forces a\vee\neg a\) for every \(U,a\); by Heyting interpretation this
says \(A\vee\neg A=U\). Also \(A\wedge\neg A=0\) by residuation. Hence
every subobject is complemented, so the topos is Boolean. Conversely if
every subobject is complemented, its Heyting negation is that complement:
if \(B\) complements \(A\), then \(B\subseteq\neg A\) by residuation,
and distributivity gives \(\neg A=(\neg A)\wedge(A\vee B)\subseteq B\).
Thus \(A\vee\neg A=U\) for every characteristic map, which by the universal
clause proves the internal formula. This supplies both directions, including
the direction left as an exercise in the source.

Assume each exponentiation functor preserves epis. To prove the internal
choice formula, by the all-arrow universal and implication clauses it
suffices, at every \(r:U\to Y^X\) internally surjective, to produce an
epimorphic refinement with a right-inverse function. The first part makes
\(h:U\times X\to U\times Y\) epic. Therefore \(h^Y\) is epic. Pull it back
along \(j:U\to(U\times Y)^Y\), the transpose of \(1_{U\times Y}\), obtaining
an epi \(e:V\to U\) and \(q:V\to(U\times X)^Y\) with \(h^Yq=je\).
Untranspose to \(\widetilde q:V\times Y\to U\times X\).
The equation becomes
\(h\widetilde q=\langle e\pi_1,\pi_2\rangle\).
Its first coordinate implies \(\pi_1\widetilde q=e\pi_1\).
Let \(s:V\to X^Y\) transpose \(\pi_2\widetilde q\); its untranspose is
\(\widetilde s\). Then \(\widetilde q=\langle e\pi_1,\widetilde s\rangle\),
and the second coordinate gives
\(\widetilde r\langle e\pi_1,\widetilde s\rangle=\pi_2\).
This is precisely the generic-element equality
\(V\times Y\forces r(e)(s(y))=y\).
Universal and existential clauses give the required internal right inverse
on the epimorphic refinement \(V\), and hence validity of choice.

Conversely suppose that formula is valid for every \(X,Y\).
Let \(f:X\to Y\) be epic and fix \(E\); we must prove \(f^E:X^E\to Y^E\)
epic. Set \(U=Y^E\), and let \(r:U\to Y^X\) be the constant family given
by the transpose of \(f\). Its untranspose is \(f\pi_2\), and
\(1_U\times f\) is epic by pullback-stability. Thus \(r\) is internally
surjective. Choice gives an epi \(e:V\to U\) and \(s:V\to X^Y\) whose
untranspose satisfies \(f\widetilde s=\pi_2\) on \(V\times Y\).
Compose \(\widetilde s\) with
\(\langle\pi_1,\operatorname{ev}_Y(e\times1_E)\rangle:V\times E\to V\times Y\)
to obtain \(t:V\times E\to X\).
Then \(ft=\operatorname{ev}_Y(e\times1_E)\).
Transposition gives \(f^E\Lambda t=e\). Since \(e\) is epic, \(f^E\) is
epic. This proves the converse for all \(E\). The produced \(s\) need only
exist on \(V\); no global section of \(f\) follows without a further choice
or projectivity premise. These arguments reconstruct the source's
VI\S6 pp.~310--315, including the continuation through pp.~313--314.
\end{proof}
\end{annexitem}

\begin{annexitem}{annex:p10.mlm.7}{thm:p10-ls1980-leibniz-equality}{Source term interpretation and Leibniz equality without circularity}
\begin{proof}
Use the named elementary-topos structural package, finite limits,
exponentials and classifier. For a finite declared context \(G\), variables
are projections, \(*\) is its terminal map, pairing is categorical pairing,
membership is evaluation, and comprehension transposes the characteristic
predicate. Propositional operations use the internal Heyting algebra on
\(\Omega\), constructed naturally on all parameters from external Heyting
operations and classification. Quantifiers classify the external projection
adjoints; their representing correspondence identifies them with internal
power-object quantifiers. Thus each source nonarithmetic constructor has
its stated arrow type, including nested formula-valued terms. Applied named
arrows use composition; optional \(0,S\) use the specified NNO arrows.

Substitution follows by induction on that grammar. Projections, unit,
pairing, composition and membership commute with substitution of the
parameter tuple. Comprehension commutes by naturality of transposition;
Heyting operations by inverse-image preservation. Alpha-rename a bound
quantifier free for the substitution; the parameter/projection square is a
pullback, and existential or universal Beck--Chevalley gives the extension
identity. Optional arithmetic constructors commute with composition too.
This supplies a generic-topos source-grammar substitution lemma, not an
appeal to the earlier presheaf-only theorem. Classifier and exponential
inverse identities make membership in comprehension exactly the original
predicate in every parameter context. Semantic variable projections are
not new projection terms added to the pure source grammar.

Let \(a,b:G\to A\). Let \(D\hookrightarrow G\) be their equalizer and
\(L\hookrightarrow G\) the extension of source Leibniz equality.
On \(D\), the arrows agree. On \(D\times PA\), their membership maps
therefore agree, so both implications between their extensions are top.
Universally quantify to get \(D\subseteq L\).
Conversely restrict to \(L\), and construct the categorical singleton of
\(a|_L\) as the graph \(K\hookrightarrow L\times A\) using the diagonal of
\(A\). Let \(k:L\to PA\) classify this graph. This helper uses the
categorical diagonal, not the source equality whose interpretation we are
proving. Universal forcing of the Leibniz formula tests every generalized
predicate, hence tests \(k\). Membership of \(a|_L\) in \(k\) holds by the
graph definition, so membership of \(b|_L\) does too. The graph pullback
says precisely \(b|_L=a|_L\). Thus \(L\subseteq D\), proving equality of
extensions, not merely of global points, also when types are empty.

The source's \S5 equality-based interpretation of its primitive logical
constructors matches the Heyting/adjoint interpretation used here.
Equality of \(\langle p,q\rangle\) with \(\langle\top,\top\rangle\) holds
exactly where both characteristic maps are truth, so classifies their meet.
Equality \((p\wedge q)=p\), at a parameter subobject \(T\hookrightarrow G\),
means equality of their restricted characteristic maps, hence
\(T\cap[p]\cap[q]=T\cap[p]\). This is \(T\cap[p]\subseteq[q]\), which
by residuation is precisely \(T\subseteq[p\Rightarrow q]\).
Finally equality of the comprehension of \(\phi\) with the full predicate
on \(A\) means its classified subobject is all of the parameter product
with \(A\); by the universal adjunction this is exactly universal
quantification of \(\phi\). Thus the source's printed \S5 interpretations
of conjunction, implication and universal quantification have the same
extensions, not merely a shared informal name. Its paragraph does not print
the remaining logical definitions; the independently validated source
\S1 adjunction rows determine the bottom, join and existential operations
used in this reconstruction.

For finite \(X\), subobject inclusion implies the assignment entailment by
pullback. Conversely test entailment at the domain \([p]\), assigned by its
inclusion into the context product, to obtain \([p]\subseteq[q]\).
For arbitrary declared \(X\), every term has finite support and can be
evaluated on that finite tuple of the assignment family. Composition with
refinements respects evaluation; no infinite product or selected inhabitants
of unused types are required. This explicitly replaces the source's
indeterminate-arrow notation in \S5 pp.~240--242 by contextual maps.
\end{proof}
\end{annexitem}

\begin{annexitem}{annex:p10.mlm.8}{thm:p10-ls1980-hol-soundness}{Every source rule direction and derivation induction}
\begin{proof}
Validity is the generalized-assignment test of
Definition~\ref{def:p10-ls1980-contextual-validity}. The already proved forcing
clauses and substitution lemma are used below, not a soundness premise.

\paragraph{Four structural rows.}
Reflexivity and transitivity are implication and composition of assignment
tests. For weakening, restrict any \(Y\)-assignment to \(X\); the two formula
values are unchanged. For substitution, evaluate \(b\) under the original
\(X\)-assignment, then form the premise assignment on \(X\cup\{y\}\) by
retaining all values except \(y\), which receives that value of \(b\).
If \(y\in X\), its original assigned value is used when evaluating \(b\),
not recomputed after replacement. The substitution lemma identifies the
premise formulas with the conclusion formulas. These arguments cover
infinite source indices because term support is finite and the extended
family is explicitly defined.

\paragraph{All ten propositional directions, logical (1)--(5).}
Top always holds. A bottom-forcing domain is initial, so factors through
every predicate, giving logical (2). Meet introduction combines both
factorizations, while its two eliminations project them.
For join elimination, forcing \(p\vee q\) gives refining domains with epic
copairing; compose the entire assignment with each refinement. Apply the
respective premises to obtain \(r\) on both domains, then descend along
the copairing to \(U\). The two reverse join rows hold because either
disjunct factors through the join.
For residual introduction, start with \(p\) at \(U\). At every refinement
\(V\to U\) forcing \(q\), persistence gives \(p\); the premise
\(p\wedge q\models_Xr\) gives \(r\). This is the implication clause.
Conversely under \(p\wedge q\), premise validity gives \(q\Rightarrow r\),
whose identity-refinement test and \(q\) give \(r\). Thus both residual
directions are validated. For finite contexts these are exactly the lattice
and residual universal properties, never Boolean material implication.

\paragraph{All four quantified directions, logical (6)--(7).}
For universal introduction, assign \(X\) at \(U\) and suppose \(p\) holds.
At every refinement \(e:V\to U\) and value \(b:V\to B\), compose the
assignment with \(e\) and assign \(y=b\), overriding its old value if it was
already declared in \(X\). The side condition \(y\notin\mathrm{FV}(p)\) makes
\(p\) unchanged; it also persists. Premise validity gives
\(\psi(b)\), hence the all-arrow universal clause gives the conclusion.
For its reverse row, restrict an \(X\cup\{y\}\)-assignment to \(X\), obtain
\(\forall y\,\psi\) from the premise and \(p\), then test at identity and
the already assigned \(y\).
For existential elimination, an assignment forcing \(\exists y\,\psi\)
gives an epi \(e:V\to U\) and witness \(b:V\to B\). Assign \(y=b\) in the
composed family, again overriding any old declared value, apply the premise
to get \(p\), then descend it to \(U\). Absence of \(y\) from \(\mathrm{FV}(p)\)
removes witness dependence. In the reverse row,
\(\psi\) under an \(X\cup\{y\}\)-assignment gives existence by identity
refinement and its assigned \(y\); the premise gives \(p\) under the
restricted \(X\)-assignment. For finite contexts these are the two
projection adjunctions, with no inhabitation assumption.

\paragraph{Six nonarithmetic nonlogical rows.}
Comprehension is the classifier/evaluation inverse identity proved in the
preceding item; its universal closure is valid.
For power extensionality, on any domain where its antecedent holds,
generic-element universal forcing gives both containments of the
subobjects of \(U\times A\) classified by \(u,v:U\to PA\). They agree;
classification and transposition give \(u=v\) as arrows. The Leibniz bridge
gives the displayed equality formula at that domain, proving its implication
at every refinement. For proposition extensionality, \(s\Leftrightarrow t\)
makes the subobjects classified by \(s,t:U\to\Omega\) contained in each
other. Uniqueness of characteristic maps makes the arrows equal; the bridge
again gives source equality.
The unit axiom holds because arrows into \(1\) agree with \(*\).
Product decomposition has witnesses \(\pi_Az,\pi_Bz\) for each generalized
\(z:U\to A\times B\). Equality of two paired arrows implies equality of
their coordinates by projections, validating product injectivity. The bridge
translates each arrow equality back to source Leibniz equality. Universal
closures and implication guards follow by the all-arrow tests. These proofs
use semantic projections, not fabricated projection terms in pure \(\mathcal L_1\).

\paragraph{Derivation induction.}
The least relation generated by these finitary rows consists of finite
derivation trees. Induct on the tree: leaves are the validated axiom rows,
and every inference preserves validity by the preceding checks. This proves
the theorem for every declared \(X\). Extra theory axioms require their
own contextual validity, not automatic acceptance. Finite hypothesis lists
reduce to the source's one-antecedent judgment by conjunction, giving the
usual contextual subobject consequence. No first-order LJ theorem has stood
in for the power and extensionality checks.

\paragraph{Optional arithmetic: derive Peano from an NNO.}
An NNO has arrows \(0:1\to N\), \(S:N\to N\) and unique recursion to each
\(a:1\to A\), \(d:A\to A\), with \(k0=a\), \(kS=dk\).
First derive recursion with parameters. Given \(a:G\to A\) and
\(d:G\times A\to A\), define \(D:A^G\to A^G\) by transposing
\((r,g)\mapsto d(g,r(g))\). Transpose \(a\) to \(\bar a:1\to A^G\),
recurse to \(k:N\to A^G\), then untranspose, exchanging product factors,
to \(h:G\times N\to A\). It satisfies
\(h(g,0)=a(g)\), \(h(g,Sn)=d(g,h(g,n))\); transposition proves uniqueness.
If \(A\to G\) has a zero section and an endomorphism over \(G\), use that
section and endomorphism in this construction. The composite \(G\times N\to G\)
is projection by recursion uniqueness. Thus \(G\times N\) is an NNO in
the slice over \(G\), a result derived here rather than presumed.

Recurse to \(c:N\to\Omega\), with \(c0=\mathsf{true}\) and step constantly
\(\mathsf{false}\). If \(Sx=0\) on \(U\), then
\(\mathsf{false}\,!_U=\mathsf{true}\,!_U\). The classifier pullback for
\(0\hookrightarrow1\) gives \(U\to0\), and strictness makes \(U\) initial.
This proves the zero Peano implication.
For successor injectivity, recurse to \(q:N\to N\times N\) from \((0,0)\)
with step \((a,b)\mapsto(Sa,a)\). Recursion uniqueness gives
\(\pi_1q=1_N\), so \(P=\pi_2q\) satisfies \(PS=1_N\). Thus \(S\) is monic;
the Leibniz bridge proves the equality axiom.
For induction, at any \(u:U\to PN\) on which the antecedent holds, let
\(R\hookrightarrow U\times N\) be its classified predicate. Membership of
zero gives a section \(U\to R\), and the universally quantified closure
condition gives a successor endomorphism of \(R\) over \(U\). The slice NNO
gives \(h:U\times N\to R\) whose composite with inclusion is identity by
recursion uniqueness. The inclusion is monic and split epic, hence an
isomorphism. Therefore \(R=U\times N\), giving the universal conclusion.
Repeating this at every refinement validates its implication and outer
predicate quantifier. None of the three Peano conclusions was an assumed
premise stating what we wished to prove.

Presheaves have the constant NNO \(N(c)=\mathbb N\) with identity
restrictions: for natural \(a:1\to A\), \(d:A\to A\), set
\(k_c(n)=d_c^n(a_c(*))\). Naturality and uniqueness follow by induction on
\(n\). Its zero and successor are componentwise ordinary zero and successor.
This supplies the actual optional presheaf arithmetic model without Choice.
A general elementary topos without an NNO gets only the nonarithmetic
fragment; the three arithmetic leaves are not unconditionally valid.

\paragraph{Optional Boolean row and topological interpretation.}
Theorem~\ref{thm:p10-mlm-internal-epi-boolean-choice} validates the Boolean
row exactly when the topos is Boolean. Under that hypothesis its leaf is
valid and the same tree induction applies; the body two-stage example
refutes unconditional validity. In sheaf models, finite-context validity
is subsheaf inclusion, so Annex~\ref{annex:p10.mlm.5} gives exactly the same
assignment tests. Join and existence then use cover-local disjuncts and
witnesses, not a pointwise disjunct or global witness. Associated-sheaf
construction and the generic-topos structural package remain explicitly
imported dependencies, not conclusions manufactured by soundness.
\end{proof}
\end{annexitem}

\begin{annexitem}{annex:p10.mlm.9}{thm:p10-ls1980-exponential-graph-bridge}{Total single-valued relations represent functions}
\begin{proof}
Comprehension and the Leibniz--diagonal bridge interpret the stated source
formula as \(\operatorname{Fun}(A,B)\hookrightarrow P(A\times B)\).
For \(k:U\to\operatorname{Fun}(A,B)\), classification gives
\(R\hookrightarrow U\times A\times B\). Its projection \(q:R\to U\times A\)
is epic: totality, tested at the generic element \(U\times A\), says its
existential image is the whole target. No preimage function is selected.
It is monic too. If \(v,w:V\to R\) have \(qv=qw\), their \(U\)- and
\(A\)-coordinates agree. Their \(B\)-coordinates both satisfy relation
membership; the single-valued clause equates them at this generalized
element. All three coordinates agree, so monicity of the relation inclusion
gives \(v=w\). Thus \(q\) is a monic epi and an isomorphism. Compose its
inverse with the \(B\)-coordinate to obtain \(h:U\times A\to B\); the
relation is precisely its graph.

Conversely \(h:U\times A\to B\) gives the graph monomorphism
\((1_{U\times A},h)\). It is total and single-valued on every generalized
context, so its classifying map factors through \(\operatorname{Fun}(A,B)\).
The constructions are inverse and commute with a parameter map \(a:U'\to U\):
pullback of a graph gives the graph of \(h(a\times1_A)\). Hence
\[
 \mathcal E(U,\operatorname{Fun}(A,B))\cong
 \mathcal E(U\times A,B)\cong\mathcal E(U,B^A)
\]
naturally in every \(U\). Apply the inverse procedures to the two identity
maps to obtain inverse arrows between the representing objects. This proves
the canonical isomorphism using all parameters, not global points or Choice.

Graph membership of \((a,b)\) is equality \(b=h(a)\) by its pullback.
For \(t:G\times A\to B\), the source comprehension in the statement has
exactly its graph: its witnesses enforce \(z=\langle x,y\rangle\) and
\(y=t(x)\); every graph pair supplies those witnesses. Local witnesses give
the same graph by epimorphic descent. Thus it represents exactly
\(\Lambda t:G\to B^A\). Evaluation gives
\(\operatorname{ev}(\Lambda t\times1_A)=t\), the beta identity.
Conversely the graph of the evaluation of \(r:G\to B^A\) recovers that same
function; transposition gives \(r\), the eta identity.
When \(A=0\), the empty relation is the unique total functional relation
and \(B^0=1\). When \(B=0\), totality permits exactly contexts with
\(U\times A\) initial, as does an arrow into \(0\). No inhabitant is used.

Pure \(\mathcal L_1\) has the \(P(A\times B)\)-valued graph comprehension,
not an implicit cast into a primitive function type. In its declared
applied extension, \(\operatorname{Fun}(A,B)\), \(B^A\), the isomorphism and
evaluation may be named objects and arrows. For an exponential expression,
its interpreted finite-context arrow is expressible as such a named
operation on the context tuple; the construction above proves what that
arrow represents. Recursive expansion of exponential macros via their
graphs and these named arrows preserves interpretation. Beta/eta replacement
therefore preserves formula characteristic maps and extensions, including
under binders by substitution. The logical validation uses those actual
interpretations, not identification of distinct source grammars. This
establishes the semantic/definitional bridge, without claiming pure-language
syntactic conservativity or a free-topos universal property.
\end{proof}
\end{annexitem}
